\SourceLeaf{217}%
\bigskip\Needspace{12\baselineskip}
\pdfbookmark[1]{Fortgesetzte Nachrichten über einen neuen von Dr. Olbers entdeckten Planeten. Juni 1802}{pallas180206}
\fancyhead[L]{\small PALLAS}
\begin{center}
{\small FORTGESETZTE NACHRICHTEN ÜBER EINEN NEUEN, HÖCHST BEMERKENSWERTEN PLANETEN UNSERES SONNENSYSTEMS, DEN DR. OLBERS IN BREMEN ENTDECKT HAT UND DER SICH ZWISCHEN MARS UND JUPITER BEWEGT.\par}
\medskip\hrule\medskip
Monatliche Correspondenz. Band V. S.\ 591—606. Juni 1802.\par
\medskip\hrule
\end{center}\medskip
Als wir im vorigen Heft zum ersten Mal von der bemerkenswerten Entdeckung des neuen Gestirns durch \textsc{Olbers} berichteten, berechtigten uns die damals vorliegenden Beobachtungen, Berechnungen und gesammelten Erfahrungen noch nicht dazu, diesen sonderbaren Wandelstern als einen Planeten unseres Sonnensystems bekannt zu geben. Wir hatten allerdings schon die Vermutung geäußert, dass dieses durch seine Bewegung und sein Aussehen so auffällige Gestirn nicht zur Klasse der gewöhnlichen sogenannten Kometen zu rechnen sei. Doch schon beim ersten Anblick dieses Himmelskörpers war die Vermutung seiner \textit{planetarischen Natur} so stark, dass sein Entdecker es sogleich wagte, für seine Bahn einen zur Sonne konzentrischen Kreis anzusetzen, ohne es wie sonst üblich zuerst mit einer Parabel zu versuchen. Aber wie wir bereits auf S.\ 486 des vorigen Heftes mitgeteilt haben, ließen sich die vorhandenen Beobachtungen durch keinerlei Kreisbahn darstellen. So wenige es auch waren, so machte es die starke Bewegung dieses Gestirns in der Breite doch schon möglich, den Kreis bei künftigen Untersuchungen dieser Bahn für immer auszuschließen.

Dasselbe versuchte Dr.\ \textsc{Gauß}. Auch er bestätigte die Erfahrung von Dr.\ \textsc{Olbers}, dass die Bewegung der \textit{Pallas} stets schneller war, als es nach den Gesetzen von \textsc{Kepler} auf einem Kreis möglich wäre. Sobald Dr.\ \textsc{Gauß} aber die Beobachtung dieses Gestirns vom 17. April von Dr.\ \textsc{Olbers} erhalten hatte, machte er sich daran, den Kegelschnitt unabhängig von jeder Hypothese nach seiner eigenen Methode zu bestimmen, deren Vorzüge \textit{Ceres} schon so glänzend bestätigt hatte. Es war zwar abzusehen, dass die Beobachtungen mit dem Ringmikrometer für diese heikle Rechnung zu ungenau und von zu kurzer Dauer waren, zumal sich \textit{Pallas} Anfang April in einer ungünstigen Lage befand, da die Tangente an ihre geozentrische Bewegung durch den Ort der Sonne ging. Dennoch fand Dr.\ \textsc{Gauß} eine Ellipse, die sich in ihren Hauptabmessungen von der später bestimmten kaum unterschied. Das Ergebnis stimmte mit dem späteren im Großen und Ganzen darin überein, dass Pallas ein Planet zwischen Mars und Jupiter sei, dessen Bahn der Ceresbahn an einer Stelle sehr nahe kommt und bei einer sehr großen Neigung eine beträchtliche Exzentrizität aufweist.

Die Exzentrizität betrug dabei $0.3$. Sie war also größer als bei allen anderen Planeten, aber doch nicht groß genug, um \textit{Pallas} den Namen eines Planeten abzusprechen. Denn das Verhältnis der großen zur kleinen Achse beträgt $1:\BareRoot{(1-(0.3)^2)}$, also ungefähr $1:0.95$. Die Ellipse ist folglich weniger abgeflacht als die Jupiterscheibe.

Da die hierfür verwendeten Beobachtungen naturgemäß bei Weitem nicht so genau sind wie Meridianbeobachtungen, misstraute Dr.\ \textsc{Gauß} diesen Elementen aus guten Gründen und hielt ihre Veröffentlichung bis zu einer späteren Prüfung zurück. Erst als er unsere \textit{Seeberger} Meridianbeobachtungen dieses Gestirns vom 15., 18. und 19. April und danach drei Pariser Beobachtungen vom 10., 12. und 13. April erhielt, entschloss er sich mangels einer längeren Reihe genauer Beobachtungen, diese vor%
\SourceLeaf{218}%
läufigen Elemente so zu verbessern, dass sie zu sechs \textit{Seeberger} Beobachtungen vom 4., 5., 7., 15., 18. und 19. April und zu den drei \textit{Pariser} Beobachtungen passten. Dabei fand er gleich beim ersten Versuch die folgenden
\expandafter\def\csname HistoricalNoteText19\endcsname{Genau diese Voraussetzung hatte Dr.\ \textsc{Gauß} auch bei \textit{Ceres} zugrunde gelegt. Er betrachtete das Aphel und den $\AscendingNode$ als ruhend gegenüber den Fixsternen und nicht gegenüber unserem Äquinoktium.}
\Needspace{20\baselineskip}
\begin{center}
\textit{Elemente der Pallas} (I).\par\medskip
[Handschrift:] Heliozentrischer Bogen $4^0\ 11'$\par\medskip
Bahn einer rechtläufigen Ellipse.\par\medskip
{\small\setlength{\tabcolsep}{2pt}
\begin{tabular}{l r}
Epoche 1802, 31. März, Mittag auf dem Seeberg $\ .\ .\ .\ .\ .\ .\ .$&$166^0\quad1'\ 37''2$\\
Tägliche mittlere tropische Bewegung $\ .\ .\ .\ .\ .\ .\ .$&$800''770$\\
Umlaufszeit $\ .\ .\ .\ .\ .\ .\ .\ .\ .\ .\ .\ .\ .\ .\ .\ .\ .$&$1618\tfrac12$ Tage\\
Logarithmus der großen Halbachse $\ .\ .\ .\ .\ .\ .\ .$&$0.4310494$
\end{tabular}
\par\smallskip
$\left.\begin{array}{l}\text{Sonnenferne}\\\text{Knoten}\end{array}\right\}\quad
\text{\shortstack{zur Zeit der Epoche\\gegenüber den Fixsternen ruhend\footnotemark}}\left\{\begin{array}{lr}.&304^0\ 36'\ 30''\\.&172^0\quad9'\ 58''\end{array}\right.$
\footnotetext{\csname HistoricalNoteText19\endcsname}
\par\smallskip
\begin{tabular}{l r}
Exzentrizität $\ .\ .\ .\ .\ .\ .\ .\ .\ .\ .\ .\ .\ .\ .\ .\ .$&$0.215708$\\
Neigung $\ .\ .\ .\ .\ .\ .\ .\ .\ .\ .\ .\ .\ .\ .\ .\ .\ .$&$33^0\ 39'\ 16''6$
\end{tabular}}
\end{center}
\Needspace{25\baselineskip}
Mit diesen Elementen verglich Dr.\ \textsc{Gauß} nun sehr genau die \textit{Seeberger} Beobachtungen, die ihm damals vorlagen, und die übrigen, die er von uns erhielt, ebenso wie die \textit{Pariser} Beobachtungen und eine von \textsc{Olbers}. Der Vergleich ergab die folgende Übereinstimmung.
\begin{center}
{\small\setlength{\tabcolsep}{3pt}
\begin{tabular}{r r|r|r|r|r|l}
\multicolumn{2}{c|}{}&&&\multicolumn{2}{c|}{Unterschiede}&\\
\multicolumn{2}{c|}{1802}&\shortstack{Berechnete\\R.A. der Pallas}&\shortstack{Berechnete\\Deklination\\der Pallas}&in R.A.&in Dekl.&Beobachter\\\hline
April&4&$183^0\ 44'\quad6''3$&$13^0\ 54'\ 51''0$&$-0''3$&$-2''0$&von Zach\\
&5&$183\ 34\ 25.7$&$14\ 13\ 17.2$&$+2.0$&$-5.7$&—\\
&7&$183\ 15\ 40.0$&$14\ 48\ 55.5$&$+1.5$&$-6.6$&—\\
&8&$183\quad6\ 35.2$&$15\quad6\quad9.1$&$-2.6$&$-0.9$&—\\
&10&$182\ 48\ 56.5$&$15\ 39\ 40.9$&$-1.5$&$-1.1$&Dr.\ Burckhardt\\
&12&$182\ 32\ 34.3$&$16\ 11\quad4.2$&$+2.8$&$+6.2$&Méchain\\
&13&$182\ 24\ 47.2$&$16\ 26\quad6.6$&$+2.8$&$+0.6$&—\\
&15&$182\ 10\ 13.2$&$16\ 54\ 33.0$&$-3.3$&$+2.2$&von Zach\\
&18&$181\ 50\ 26.2$&$17\ 34\ 26.5$&$-4.4$&$\ .\ .\ .$&—\\
&19&$181\ 44\ 28.0$&$17\ 46\ 51.7$&$+2.7$&$-2.5$&—\\
&24&$181\ 19\ 39.4$&$18\ 42\ 27.8$&$+1.6$&$\ .\ .\ .$&—\\
&25&$181\ 15\ 44.1$&$18\ 52\ 18.3$&$+11.9$&$\ .\ .\ .$&—\\
&26&$181\ 12\ 10.4$&$19\quad1\ 43.5$&$+8.6$&$-5.9$&—\\
&27&$181\quad8\ 57.9$&$19\ 10\ 44.0$&$+12.2$&$-2.8$&—\\
&29&$181\quad3\ 39.6$&$19\ 27\ 31.2$&$+23.0$&$-12.5$&—\\
&30&$181\quad1\ 33.9$&$19\ 35\ 18.2$&$+24.0$&$-15.3$&—\\
Mai&1&$180\ 59\ 48.1$&$19\ 42\ 42.5$&$+25.6$&$+1.7$&—\\
&1&$180\ 59\ 38.0$&$19\ 43\ 38.0$&$+20.0$&$+7.0$&Dr.\ Olbers
\end{tabular}}
\end{center}
Addiert man die Unterschiede mit umgekehrten Vorzeichen zu den berechneten Werten, erhält man die Beobachtungen.

Auch hier hält Dr.\ \textsc{Gauß} wie bei \textit{Ceres} an seiner bisherigen, nachahmenswerten Gewohnheit fest, die Rektaszensionen und Deklinationen aus den Elementen selbst abzuleiten und mit den beobachteten zu vergleichen, ohne aus den beobachteten Rektaszensionen und Deklinationen \textit{beobachtete Längen und Breiten abzuleiten}. Der Grund dafür leuchtet so unmittelbar ein, dass man sich wundern muss, warum
\SourceLeaf{219}%
dieses Verfahren nicht schon längst allgemein eingeführt und angewandt worden ist. Der Aufwand ist gleich, ob man mit der Beobachtung der Rechnung einen Schritt entgegengeht oder die Rechnung ganz bis zur Beobachtung führt. Nach einem eigenen Verfahren von Dr.\ \textsc{Gauß} wird diese Parallaxenrechnung vielmehr noch etwas einfachen. Der größte Vorteil liegt aber darin, dass man die guten Beobachtungen von den möglicherweise weniger zuverlässigen trennen kann. Andernfalls kann eine falsche Deklination leicht eine gute Rektaszension verderben, sodass sich sowohl eine fehlerhafte Breite als auch eine schlechte Länge ergibt. Da im obigen Vergleich der Beobachtungen mit den Bahnelementen die Unterschiede gegen Ende immer größer werden, war abzusehen, dass ein Unterschied von $25''$ bei der kurzen Dauer der Beobachtungen noch einen sehr starken Einfluss auf die Bestimmung dieser Elemente haben musste. Die weitere Untersuchung, die Dr.\ \textsc{Gauß} unmittelbar nach diesem Vergleich vornahm, hatte ihm auch gezeigt, dass die Beobachtungen über 15 Tage, auf denen seine früheren Elemente beruhten, noch keine genaue Bahnbestimmung erlaubten, die Natur des Gestirns aber schon recht deutlich erkennen ließen. Er hat nun sofort aus der ganzen Reihe unserer \textit{Seeberger} Beobachtungen neue Elemente bestimmt, die möglichst gut zu diesen Beobachtungen vom 4. April bis zum 1. Mai passen. Da der auf der Bahn durchlaufene Bogen bereits $7^0\ 26'$ beträgt, dürfte man zumindest an der planetarischen Natur des von \textsc{Olbers} entdeckten Gestirns wohl nicht mehr zweifeln. Die neue Ellipse dieses planetarischen Himmelskörpers ist somit die folgende.
\begin{center}
\textit{Elemente der Pallas aus 27 Tagen Seeberger Beobachtungen.}\par\medskip
(II.)\par\medskip
{\small\setlength{\tabcolsep}{2pt}
\begin{tabular}{l r}
Epoche 1802, 31. März, Mittag auf dem Seeberg $\ .\ .\ .\ .\ .\ .\ .$&$161^0\ 12'\ 43''2$\\
Tägliche mittlere tropische Bewegung $\ .\ .\ .\ .\ .\ .\ .\ .$&$757''166$\\
Logarithmus der großen Halbachse $\ .\ .\ .\ .\ .\ .\ .\ .$&$0.4472636$
\end{tabular}
\par\smallskip
$\left.\begin{array}{l}\text{Sonnenferne}\\\text{Knoten}\end{array}\right\}\ \text{gegenüber den Fixsternen ruhend}\left\{\begin{array}{lr}
.\ .\ .\ .\ .\ .\ .\ .\ .&300^0\quad5'\quad4''\\
.\ .\ .\ .\ .\ .\ .\ .\ .&172\ 34\ 35
\end{array}\right.$
\par\smallskip
\begin{tabular}{l r}
Exzentrizität $\ .\ .\ .\ .\ .\ .\ .\ .\ .\ .\ .\ .\ .\ .\ .\ .$&$0.2591096$\\
Neigung $\ .\ .\ .\ .\ .\ .\ .\ .\ .\ .\ .\ .\ .\ .\ .\ .\ .\ .$&$35^0\quad0'\ 42''$
\end{tabular}}
\end{center}

Dr.\ \textsc{Gauß} verglich nun die ganze Reihe unserer \textit{Seeberger} Beobachtungen mit diesen Elementen und fand die folgende Übereinstimmung.
\begin{center}
{\small\setlength{\tabcolsep}{5pt}
\begin{tabular}{r r|r|r|r|r}
\multicolumn{2}{c|}{}&&&\multicolumn{2}{c}{Unterschiede}\\
\multicolumn{2}{c|}{1802}&\shortstack{Berechnete\\R.A.\\der Pallas}&\shortstack{Berechnete nördliche\\Deklination\\der Pallas}&in R.A.&in Dekl.\\\hline
April&4&$183^0\ 44'\quad5''7$&$13^0\ 54'\ 55''4$&$-0''9$&$+3''4$\\
&5&$183\ 34\ 24.6$&$14\ 13\ 22.6$&$+0.9$&$-0.3$\\
&7&$183\ 15\ 39.6$&$14\ 49\quad1.3$&$+1.6$&$-0.7$\\
&8&$183\quad6\ 35.4$&$15\quad6\ 15.2$&$-2.4$&$+5.2$\\
&15&$182\ 10\ 17.9$&$16\ 54\ 35.7$&$+1.4$&$+4.9$\\
&18&$181\ 50\ 29.4$&$17\ 34\ 27.1$&$-1.2$&$\ .\ .\ .$\\
&19&$181\ 44\ 30.7$&$17\ 46\ 51.9$&$+5.4$&$-2.5$\\
&24&$181\ 19\ 35.8$&$18\ 42\ 27.9$&$-2.0$&$\ .\ .\ .$\\
&25&$181\ 15\ 38.3$&$18\ 52\ 18.9$&$+6.1$&$\ .\ .\ .$\\
&26&$181\ 12\quad2.0$&$19\quad1\ 44.9$&$+0.2$&$-4.5$\\
&27&$181\quad8\ 46.6$&$19\ 10\ 46.3$&$+1.0$&$-0.5$
\end{tabular}
\par\medskip
\begin{tabular}{r r|r|r|r|r}
\multicolumn{2}{c|}{}&&&\multicolumn{2}{c}{Unterschiede}\\
\multicolumn{2}{c|}{1802}&\shortstack{Berechnete\\R.A.\\der Pallas}&\shortstack{Berechnete nördliche\\Deklination\\der Pallas}&in R.A.&in Dekl.\\\hline
April&29&$181^0\quad3'\ 21''2$&$19^0\ 27'\ 36''2$&$+4''6$&$-7''5$\\
&30&$181\quad1\ 11.4$&$19\ 35\ 24.9$&$+1.5$&$-8.6$\\
Mai&1&$180\ 59\ 22.9$&$19\ 42\ 50.8$&$+0.4$&$+10.0$\\
&2&$180\ 57\ 56.7$&$19\ 49\ 53.2$&$+4.8$&$+1.6$\\
&3&$180\ 56\ 52.4$&$19\ 56\ 32.9$&$-3.2$&$-14.8$\\
&5&$180\ 55\ 49.6$&$20\quad8\ 45.6$&$-5.0$&$+6.9$\\
&6&$180\ 55\ 51.4$&$20\ 14\ 19.4$&$-2.7$&$\ .\ .\ .$\\
&7&$180\ 56\ 14.9$&$20\ 19\ 32.1$&$-8.9$&$-3.2$\\
&8&$180\ 57\quad0.4$&$20\ 24\ 23.7$&$-9.6$&$+3.6$\\
&11&$181\quad1\ 25.5$&$20\ 36\ 58.8$&$-16.2$&$-14.4$
\end{tabular}}
\end{center}
\SourceLeaf{220}%

Dr.\ \textsc{Gauß} hätte sich sehr leicht noch etwas genauer an die letzten Rektaszensionen anschließen können. Wegen der bei herannahender Dämmerung weniger zuverlässigen Deklinationen hätte er aber nicht sicher sein können, ob sich dadurch wirklich eine \textit{Verbesserung} der Bahn ergäbe. Er will daher spätere Beobachtungen abwarten, bevor er diese Elemente erneut korrigiert.

Wie dem auch sei, alle bisherigen Beobachtungen, Rechnungen und Vergleiche scheinen doch die planetarische Natur der \textit{Pallas} zu begründen und uns vor allen Zweifeln, Bedenken und unbestimmten Hypothesen zu bewahren. Dr.\ \textsc{Gauß} bemerkt nämlich in einem Schreiben ganz treffend: „Einer genauen Rechnung kann man nur eine ebenso genaue Rechnung entgegensetzen. Leere Vermutungen und vage Überlegungen genügen dazu nicht. Um Pallas im Jahr 1803 leicht wiederzufinden, werden meine Elemente II allerdings noch beträchtliche Verbesserungen brauchen und sie hoffentlich auch erhalten, wenn nur noch genügend Beobachtungen dafür gelingen. Doch \textit{beträchtlich} werden diese Verbesserungen nur für die Bestimmung des Ortes im Jahr 1803 sein. Für die \textit{planetarische} oder \textit{nichtplanetarische Natur} der \textit{Pallas} sind sie es nicht.“

Schon nach den früher gefundenen Elementen kam die Bahn der \textit{Pallas} der Ceresbahn sehr nahe, dort, wo ihr aufsteigender Knoten auf der Ceresbahn liegt. Das brachte Dr.\ \textsc{Olbers} auf die Vermutung, dass sich die Bahnen der \textit{Pallas} und der \textit{Ceres} tatsächlich schneiden und vielleicht wie zwei Glieder einer Kette ineinandergreifen. Dr.\ \textsc{Gauß} hat die Abstände der \textit{Pallas} und der \textit{Ceres} auf der Knotenlinie ihrer Bahnen überschlagen. Er fand bei $\DescendingNode$ für den Abstand der \textit{Pallas} von der Sonne $2.86$ und für den der \textit{Ceres} $2.93$. Bei $\AscendingNode$ sind die Abstände weniger ähnlich. Die Reihe der Beobachtungen ist aber noch zu kurz, um die Möglichkeit ausschließen zu können, dass die Abstände an dieser Stelle vielleicht genau gleich sind. Dr.\ \textsc{Gauß} will der Rechnung ganz ihren Lauf lassen und schlechterdings nichts Hypothetisches einmischen. Es ist freilich noch zu früh, sich in Hypothesen zu vertiefen. Wenn es gelingt, \textit{Pallas} im nächsten Jahr wiederzufinden, werden wir stattdessen bald Gewissheit haben.

Zu welchen Spekulationen über die Entstehung und Geschichte unseres Planetensystems wird dieser sonderbare Himmelskörper unterdessen Anlass geben? Welch unerwartet große Neigung für einen Planeten! Und wie auffällig ist die Lage seiner Bahn gegenüber der Ceresbahn! Diese brachte Dr.\ \textsc{Olbers} auf den Gedanken, dass beide Planeten wohl nur \textit{Überreste} eines einzigen sein könnten, den der Stoß eines Kometen zertrümmert hat. Wer hätte so etwas jemals in unserem Planetensystem erwartet? Und zu welchen neuen, wichtigen und großen Erkenntnissen kann und wird uns dieser Himmelskörper führen, der in mehr als einer Hinsicht zu den allerwichtigsten astronomischen Entdeckungen gehört, die je gemacht wurden!

Ebenso bemerkenswert findet Dr.\ \textsc{Gauß} die große Übereinstimmung der mittleren Bewegung dieser beiden Planeten. Bisher lässt sich nicht mit Gewissheit entscheiden, ob sie vielleicht \textit{genau} gleich sind. Wir kennen ja die eigentliche tägliche. Bewegung der \textit{Ceres} noch nicht auf eine Sekunde genau. So klein die Massen von \textit{Ceres} und \textit{Pallas} auch sind, muss es doch für den Unterschied ihrer mittleren Bewegungen eine Grenze geben, unter die er nicht fallen kann, ohne durch die wechselseitige Einwirkung ganz auf null gebracht zu werden. Nun unterscheiden sich diese mittleren Bewegungen nach den Elementen II nur wenig und kommen einander bei genauerer Kenntnis der Bahnen vielleicht noch viel näher. So könnte es doch sein, dass \textit{Ceres} und \textit{Pallas} dieselbe Umlaufzeit haben und nie sehr weit voneinander entfernt sind, ohne einander seit Anbeginn je zu nahe gewesen zu sein oder sich in Zukunft zu nahe kommen zu können. Bei $\AscendingNode$ wäre nämlich stets \textit{Pallas} voraus und bei $\DescendingNode$ \textit{Ceres}, weil ihre Mittelpunktsgleichung so verschieden ist. Dann wird man auch nach einer langen Reihe von Jahren aus Beobachtungen der \textit{Ceres} die
\SourceLeaf{221}%
Masse der \textit{Pallas} und aus Beobachtungen der \textit{Pallas} die Masse der \textit{Ceres} sehr gut bestimmen können. Nach einem Überschlag von Dr.\ \textsc{Gauß} ergibt ein Unterschied von 12 Sekunden in der mittleren täglichen Bewegung einen synodischen Umlauf von beinahe dreihundert Jahren.

\begingroup\rightskip=0pt plus1em\relax
Zu welchen Untersuchungen wird ein so einzigartiges Phänomen Anlass geben? Welche neuen Bedürfnisse der Wissenschaft wird es wecken? Nach allem, was wir bisher wussten, hätten wir solche außerordentlichen Erscheinungen im Weltsystem nicht erwartet! Wie sehr werden uns unsere bisherigen üblichen Rechenmethoden im Stich lassen und sich als unzureichend erweisen! Man denke etwa nur daran, dass wir die Störungen dieses neuen Himmelskörpers berechnen sollen. Nach den bisher allein üblichen Methoden hat man die Exzentrizität und die Neigung der Bahnen gleichsam als unendlich kleine Größen betrachtet und nach ihren Potenzen entwickelt, wobei man die höheren vernachlässigte. Bei zwei so großen Elementen der Pallasbahn wird das aber nicht mehr möglich sein. Ihre höheren Potenzen dürfen gewiss nicht vernachlässigt werden. Daher begründet vielleicht auch \textit{Pallas} zwangsläufig eine neue Epoche in der physischen Astronomie. Was wir durch strenge Rechnung wissen, ist ja das Einzige, was wir über den Bau der Welt \textit{wirklich wissen}. Wer rechnet, braucht aber Erfahrungen und Beobachtungen. Nur durch sie kann man Hypothesen den Stempel der Wahrheit aufdrücken. Deshalb haben alle Astronomen \textit{Pallas} mit ebenso großem Fleiß wie großer Sorgfalt beobachtet. Wir teilen unseren Lesern die ganze gesammelte Reihe dieser kostbaren Beobachtungen mit. Unsere Zeitschrift war bisher das Archiv aller Beobachtungen der \textit{Ceres}. Das soll sie auch für \textit{Pallas} sein, und wir werden keine Mühe scheuen, diese Sammlung möglichst vollständig zu machen. So finden die Astronomen hier alles beisammen, was sie sich sonst mühsam aus verstreuten Blättern zusammensuchen müssten und erst zu spät, vielleicht auch nie, finden würden.\par\endgroup

—\quad —\quad —

Um den Astronomen mit großen Äquatorialinstrumenten das Auffinden der \textit{Pallas} zu erleichtern, hat Dr.\ \textsc{Gauß} nach seinen Bahnelementen II die folgende Ephemeride ihres Laufs berechnet. Es mag etwas gewagt erscheinen, aus Elementen, die nur auf Beobachtungen aus 27 Tagen beruhen, den Ort des Planeten 2 Monate im Voraus bestimmen zu wollen. Es kann zwar nicht ausbleiben, dass diese Ephemeride gegen Ende schon sehr merklich vom Himmel abweicht. Da diese Abweichung aber wahrscheinlich nur einige Minuten betragen kann, wird sie das bequeme Auffinden des Planeten mithilfe der Ephemeride nicht erschweren, zumal man schon voraussehen kann, dass der Fehler in der Rektaszension (und vermutlich auch in der Deklination) negativ sein wird.
\Needspace{12\baselineskip}
\begin{center}
\textit{Ephemeride für Pallas Olbersiana, für Mitternacht auf dem Seeberg.}\par\medskip
{\small\setlength{\tabcolsep}{5pt}
\begin{tabular}{r r|r|r@{\qquad}r r|r|r}
\multicolumn{2}{c|}{1802}&\shortstack{R.A.\\der Pallas}&\shortstack{Deklination\\der Pallas}&\multicolumn{2}{c|}{1802}&\shortstack{R.A.\\der Pallas}&\shortstack{Deklination\\der Pallas}\\\hline
Mai&24&$181^0\ 57'$&$21^0\quad1'$ N.&Juni&14&$185^0\ 12'$&$20^0\ 17'$ N.\\
&27&$182\ 18$&$21\quad0$&&17&$185\ 48$&$20\quad5$\\
&30&$182\ 41$&$20\ 57$&&20&$186\ 27$&$19\ 52$\\
Juni&2&$183\quad6$&$20\ 52$&&23&$187\quad7$&$19\ 37$\\
&5&$183\ 34$&$20\ 46$&&26&$187\ 49$&$19\ 22$\\
&8&$184\quad5$&$20\ 38$&&29&$188\ 32$&$19\quad6$\\
&11&$184\ 37$&$20\ 28$&&&&
\end{tabular}}
\end{center}
\begin{flushright}[\textsc{von Zach}.]\end{flushright}
